% Add additional packages you would like to use here.

% %%%%%%%%%%%%%%%%%%%%%%%%%%%%%%%%%%%%%%%%%%%%%%%%%%%%%%%%%%%%%%%%%%%%%%%%%%%%%%%%%%%%%%%%%%%%
 % For professional tables
\usepackage{dcolumn}% Align table columns on decimal point
\usepackage{booktabs} % if you're using \toprule, \midrule, \bottomrule
\usepackage{tabularx}
\usepackage{multirow}
\renewcommand{\arraystretch}{1.5}
\usepackage{threeparttable}
% %%%%%%%%%%%%%%%%%%%%%%%%%%%%%%%%%%%%%%%%%%%%%%%%%%%%%%%%%%%%%%%%%%%%%%%%%%%%%%%%%%%%%%%%%%%%

% %%%%%%%%%%%%%%%%%%%%%%%%%%%%%%%%%%%%%%%%%%%%%%%%%%%%%%%%%%%%%%%%%%%%%%%%%%%%%%%%%%%%%%%%%%%%
% for chemical forumals
\usepackage{chemformula}
% %%%%%%%%%%%%%%%%%%%%%%%%%%%%%%%%%%%%%%%%%%%%%%%%%%%%%%%%%%%%%%%%%%%%%%%%%%%%%%%%%%%%%%%%%%%%


\usepackage{gensymb} % for degree ° symbol
% \usepackage{mycommands}
% \usepackage{float}
%\usepackage{floatrow}
% \usepackage{dblfloatfix}
% \usepackage{placeins}

% %%%%%%%%%%%%%%%%%%%%%%%%%%%%%%%%%%%%%%%%%%%%%%%%%%%%%%%%%%%%%%%%%%%%%%%%%%%%%%%%%%%%%%%%%%%%
% SI unit
\usepackage{siunitx}
\DeclareSIUnit\torr{torr}
\DeclareSIUnit\ppm{ppm}
\DeclareSIUnit\ppb{ppb}
\DeclareSIUnit\bar{bar}
\DeclareSIUnit\gauss{G}
\DeclareSIUnit\neV{\nano\electronvolt}
\DeclareSIUnit\peV{\pico\electronvolt}
\DeclareSIUnit\kms{km/s}
\sisetup{mode=text}
\sisetup{print-unity-mantissa=false}  % To display only the exponent (i.e., 10⁻¹¹) 
\sisetup{separate-uncertainty=false} % use parenthesis style for uncertainty
% \sisetup{uncertainty-mode=separate} % use +/- style for uncertainty
\sisetup{uncertainty-mode=compact-marker} % turn this on to use style like 123.45(1.20)
% otherwise turn this on to use style like 123.45(120)
% %%%%%%%%%%%%%%%%%%%%%%%%%%%%%%%%%%%%%%%%%%%%%%%%%%%%%%%%%%%%%%%%%%%%%%%%%%%%%%%%%%%%%%%%%%%%

% \usepackage[normalem]{ulem}
\usepackage{changes}


% \NeedsTeXFormat{LaTeX2e}
% \ProvidesPackage{mycommands}[define custom commands for easy commenting here]
\usepackage{xcolor}

\newcommand{\JW}[2][cyan]{\textcolor{#1}{Julian: #2}}

\newcommand{\OM}[2][purple]{\textcolor{#1}{Olympia: #2}}

\newcommand{\HB}[2][orange]{\textcolor{#1}{Hendrik: #2}}

\newcommand{\YZ}[1]{\textcolor{brown}{Yuzhe: #1}}

\newcommand{\kp}[2][red]{\textcolor{#1}{* #2}}

\newcommand{\new}[2][red]{\textcolor{#1}{#2}}

\newcommand{\NLF}[2][teal]{\textcolor{#1}{Nate: #2}}

\newcommand{\YG}[2][magenta]{\textcolor{#1}{Younggeun: #2}}

\newcommand{\DB}[2][blue]{\textcolor{#1}{Dima: #2}}

% \endinput
